\chapter{Tensors and the Relativistic Lorentz Force}

Tonight we turn from scalar fields to electrodynamics.  The preceding
discussion established a useful pattern: one action describes both how a
field acts on a particle and how the particle acts back on the field.  We
shall eventually apply that pattern to the electromagnetic field, but first
we need notation strong enough to expose the symmetry of the calculation.

This preliminary work is not decoration.  A minus sign, a zero, an equals
sign, and vector notation each compress an idea that would otherwise have to
be repeated in words.  Tensor notation does the same for relativity.  Once
upper and lower indices, contractions, and transformation laws are in place,
the electric and magnetic forces will assemble themselves into one
four-dimensional equation.

\section{Four-Vectors and the Economy of Indices}

\subsection{Upper and lower components}

Greek indices run over all four spacetime coordinates, while Latin indices
run only over the three spatial coordinates:
\begin{equation}
\mu,\nu,\rho,\sigma=0,1,2,3,
\qquad
m,n=1,2,3.
\end{equation}
A generic contravariant four-vector is \(A^\mu=(A^0,A^1,A^2,A^3)\).
With the convention used here,
\begin{equation}
\eta_{\mu\nu}
=
\begin{pmatrix}
-1&0&0&0\\
0&1&0&0\\
0&0&1&0\\
0&0&0&1
\end{pmatrix},
\qquad
A_\mu=\eta_{\mu\nu}A^\nu .
\label{eq:sr6-metric-lowering}
\end{equation}
Thus lowering an index changes the sign of the time component and leaves the
spatial components alone:
\begin{equation}
A_0=-A^0,
\qquad
A_m=A^m.
\label{eq:sr6-component-lowering}
\end{equation}
For flat spacetime this is all the operation amounts to, although the same
notation will acquire a richer geometric meaning in general relativity.

\begin{figure}[htbp]
\centering
\includegraphics[width=0.90\textwidth]{lecture_06_figure_01.png}
\caption{The Minkowski metric and the component rule for lowering a four-vector index.}
\lectureframe{00:07:45}
\end{figure}

Einstein's summation convention is used when an index is repeated once
upstairs and once downstairs.  The contraction
\begin{equation}
A_\mu B^\mu
=
-A^0B^0+A^1B^1+A^2B^2+A^3B^3
\label{eq:sr6-vector-contraction}
\end{equation}
is therefore a scalar.  If a sum does not have this upper--lower repeated
form, we shall display its summation sign explicitly.

\subsection{The derivative carries a lower index}

The four derivative operators are collected into
\begin{equation}
\partial_\mu
\equiv
\frac{\partial}{\partial x^\mu}.
\end{equation}
Acting on a scalar field \(\phi(x)\), they form the covariant vector
\(\partial_\mu\phi\).  The quickest way to see the index placement is to
consider an infinitesimal displacement:
\begin{equation}
d\phi
=
\frac{\partial\phi}{\partial x^\mu}\,dx^\mu
=
\partial_\mu\phi\,dx^\mu.
\label{eq:sr6-total-differential}
\end{equation}
The left-hand side is the difference of scalar values at neighboring
events, hence a scalar.  Since \(dx^\mu\) is contravariant, the object
contracted with it must be covariant.

\begin{classroomqa}
\audiencequestion{If \(x^\mu\) has an upper index, why does differentiation with respect to \(x^\mu\) produce a lower index?}
\lecturerresponse{Because the differential change \(d\phi=\partial_\mu\phi\,dx^\mu\) is a scalar.  A contravariant displacement can produce that scalar only by contracting with a covariant gradient.  The position of the derivative index is therefore fixed by its transformation law, not by the typography of the fraction alone. \lecturetimestamp{00:10:10}}
\end{classroomqa}

\begin{figure}[htbp]
\centering
\includegraphics[width=0.90\textwidth]{lecture_06_figure_02.png}
\caption{The scalar differential \(d\phi\) identifies \(\partial_\mu\phi\) as a covariant vector.}
\lectureframe{00:11:05}
\end{figure}

The same rule provides another scalar.  If \(B^\mu(x)\) is a four-vector
field, then
\begin{equation}
\partial_\mu B^\mu
\label{eq:sr6-four-divergence}
\end{equation}
is its four-divergence.  It combines the time derivative of the time
component with the three spatial derivatives of the corresponding spatial
components.

\begin{figure}[htbp]
\centering
\includegraphics[width=0.90\textwidth]{lecture_06_figure_03.png}
\caption{Two scalar contractions: \(A_\mu B^\mu\) and the four-divergence \(\partial_\mu B^\mu(x)\).}
\lectureframe{00:14:00}
\end{figure}

All of this has a three-dimensional analogue.  In Euclidean three-space the
metric is the identity matrix, so raising or lowering a Latin index changes
nothing.  We shall consequently move spatial indices up and down when it is
convenient, while continuing to distinguish spacetime indices carefully.

\section{Lorentz Transformations in Matrix Form}

Scalars and vectors are defined by the way they transform.  A scalar keeps
the same value, whereas a contravariant four-vector obeys
\begin{equation}
(A^\mu)'=L^\mu{}_\nu A^\nu.
\label{eq:sr6-vector-transform}
\end{equation}
The visible index \(\mu\) is free and occurs on both sides.  The repeated
index \(\nu\) is summed and disappears from the final component equation.

Here \emph{Lorentz transformation} includes spatial rotations as well as
boosts.  A boost may point along any axis; rotations can be composed with
boosts to produce a general transformation.  A spatial rotation itself does
not alter the time coordinate.

\begin{classroomqa}
\audiencequestion{When we rotate space, is time left completely unaffected?}
\lecturerresponse{Yes.  In an ordinary spatial rotation the time component is unchanged.  Such rotations are included in the Lorentz group, alongside transformations that mix time with space. \lecturetimestamp{00:16:27}}
\end{classroomqa}

For a boost of speed \(v\) along the \(x\)-axis, with \(c=1\),
\begin{equation}
L^\mu{}_\nu
=
\begin{pmatrix}
\gamma&-\gamma v&0&0\\
-\gamma v&\gamma&0&0\\
0&0&1&0\\
0&0&0&1
\end{pmatrix},
\qquad
\gamma=\frac{1}{\sqrt{1-v^2}}.
\label{eq:sr6-boost-matrix}
\end{equation}
Multiplication by \((t,x,y,z)^{\mathsf T}\) gives
\begin{equation}
\begin{aligned}
t'&=\gamma(t-vx),\\
x'&=\gamma(x-vt),\\
y'&=y,\\
z'&=z.
\end{aligned}
\end{equation}
The class supplies the zeros and the two unit entries as the matrix is
completed on the board.

\begin{figure}[htbp]
\centering
\includegraphics[width=0.90\textwidth]{lecture_06_figure_04.png}
\caption{The \(x\)-directed boost matrix acting on the column \((t,x,y,z)^{\mathsf T}\).}
\lectureframe{00:21:35}
\end{figure}

A rotation through angle \(\theta\) in the \(y\)--\(z\) plane instead has
\begin{equation}
R_{yz}(\theta)
=
\begin{pmatrix}
1&0&0&0\\
0&1&0&0\\
0&0&\cos\theta&\sin\theta\\
0&0&-\sin\theta&\cos\theta
\end{pmatrix}.
\label{eq:sr6-yz-rotation}
\end{equation}
It says
\begin{equation}
y'=y\cos\theta+z\sin\theta,
\qquad
z'=-y\sin\theta+z\cos\theta,
\end{equation}
while \(t\) and \(x\) remain fixed.  Products of matrices such as
\eqref{eq:sr6-boost-matrix} and \eqref{eq:sr6-yz-rotation} generate more
general Lorentz transformations.

\begin{figure}[htbp]
\centering
\includegraphics[width=0.90\textwidth]{lecture_06_figure_05.png}
\caption{A pure rotation in the \(y\)--\(z\) plane, with \(t\) and \(x\) untouched.}
\lectureframe{00:23:10}
\end{figure}

A covariant vector transforms with a matrix carrying the opposite index
placement:
\begin{equation}
(A_\mu)'=M_\mu{}^\nu A_\nu.
\label{eq:sr6-covector-transform}
\end{equation}
An audience correction at the board catches a temporarily repeated free
index: the unsummed \(\mu\) on the left must also be the unsummed index on
the right.  The two matrices represent the same physical change of frame
and are related by
\begin{equation}
M=\eta L\eta,
\qquad
\eta^{-1}=\eta.
\label{eq:sr6-M-L-relation}
\end{equation}
The last identity follows immediately because every diagonal entry of
\(\eta\) is its own inverse.

\begin{figure}[htbp]
\centering
\includegraphics[width=0.90\textwidth]{lecture_06_figure_06.png}
\caption{Contravariant and covariant transformation laws, related by \(M=\eta L\eta\).}
\lectureframe{00:25:20}
\end{figure}

The central idea is now available: a four-vector is an object whose
components transform in parallel with the spacetime coordinates.

\section{Tensors: One Transformation for Each Index}

\subsection{From vector products to the definition}

A scalar is a rank-zero tensor, a vector is a rank-one tensor, and an object
with two indices is a rank-two tensor.  Begin with two vectors and form
\begin{equation}
T^{\mu\nu}=A^\mu B^\nu.
\label{eq:sr6-simple-tensor}
\end{equation}
There are \(4\times4=16\) components. They range from \(A^0B^0\) and
\(A^0B^1\) through \(A^3B^3\).
The symbol \(T^{\mu\nu}\) denotes this entire array, not one number.

Transforming each vector separately gives
\begin{align}
(A^\mu)'&=L^\mu{}_\sigma A^\sigma,\\
(B^\nu)'&=L^\nu{}_\tau B^\tau,
\end{align}
and therefore
\begin{equation}
(T^{\mu\nu})'
=
L^\mu{}_\sigma L^\nu{}_\tau T^{\sigma\tau}.
\label{eq:sr6-rank-two-transform}
\end{equation}
Each tensor index receives its own Lorentz matrix.  The dummy indices
\(\sigma\) and \(\tau\) may be renamed at will, provided each pair remains
consistent.

\begin{classroomqa}
\audiencequestion{Does this rearrangement assume that the matrices and vector components commute?}
\lecturerresponse{These are ordinary classical numbers, not noncommuting quantum operators, so they commute.  Moreover, the same Lorentz matrix is acting on both indices.  There is no ordering issue in this calculation. \lecturetimestamp{00:31:56}}
\end{classroomqa}

The rule extends immediately to rank three:
\begin{equation}
(T^{\mu\nu\lambda})'
=
L^\mu{}_\sigma
L^\nu{}_\tau
L^\lambda{}_\kappa
T^{\sigma\tau\kappa},
\label{eq:sr6-rank-three-transform}
\end{equation}
and similarly to any rank.  This transformation law, rather than a
particular factorization into vectors, is the definition of a tensor.
For example,
\begin{equation}
A^\mu B^\nu+C^\mu D^\nu
\end{equation}
is still a tensor even though it cannot generally be written as one product
of two vectors.

\begin{figure}[htbp]
\centering
\includegraphics[width=0.90\textwidth]{lecture_06_figure_07.png}
\caption{The board ladder from a product of vectors to rank-two and rank-three tensor transformation laws.}
\lectureframe{00:33:20}
\end{figure}

\begin{classroomqa}
\audiencequestion{What does matrix multiplication by a rank-two tensor mean component by component?}
\lecturerresponse{Fix the two output indices and sum over the two hidden indices.  For example,
\[
(T^{31})'=L^3{}_\sigma L^1{}_\tau T^{\sigma\tau}.
\]
Since both \(\sigma\) and \(\tau\) take four values, this one output component contains sixteen terms. \lecturetimestamp{00:34:35}}
\end{classroomqa}

\begin{figure}[htbp]
\centering
\includegraphics[width=0.90\textwidth]{lecture_06_figure_08.png}
\caption{Expanding the fixed component \((T^{31})'\) makes the two independent summations explicit.}
\lectureframe{00:35:55}
\end{figure}

Tensor equations have a decisive advantage.  If two tensors have equal
components in one frame, the transformation law makes them equal in every
frame.  Equivalently, if every component of a tensor vanishes in one frame,
the whole tensor vanishes in every frame.  One vanishing component is not
enough; all components must vanish.  This is why an equation written between
tensors transports its truth from one inertial frame to another.

\subsection{Raising and lowering tensor indices}

An index of a tensor is lowered exactly as an index of a vector:
\begin{equation}
T^\mu{}_\nu=T^{\mu\sigma}\eta_{\sigma\nu}.
\label{eq:sr6-tensor-lowering}
\end{equation}
For the diagonal metric in \eqref{eq:sr6-metric-lowering}, every lowered time
index contributes one minus sign.  Consequently,
\begin{equation}
T_{00}=T^{00},
\qquad
T_{01}=-T^{01}.
\label{eq:sr6-tensor-component-signs}
\end{equation}
The first equality contains two sign changes; the second contains only one.

\begin{figure}[htbp]
\centering
\includegraphics[width=0.90\textwidth]{lecture_06_figure_09.png}
\caption{Examples of lowering two tensor indices: two time indices give two minus signs, while one time index gives one.}
\lectureframe{00:41:00}
\end{figure}

\begin{classroomqa}
\audiencequestion{Why is it useful to keep both covariant and contravariant forms?  Is there more to it than compact notation?}
\lecturerresponse{For the present calculation, the immediate benefit is compact contraction and automatic management of the Minkowski signs.  There is a deeper geometric distinction between vectors and covectors, but it is better developed in general relativity.  In flat spacetime the index rules are so efficient that one usually works with them directly. \lecturetimestamp{00:41:26}}
\end{classroomqa}

\begin{classroomqa}
\audiencequestion{Does a tensor always have four components, or is that peculiar to relativity?}
\lecturerresponse{The range of each index is set by the dimension of the space.  A rank-one tensor in four-dimensional spacetime has four components; a rank-two tensor has sixteen.  In three dimensions the corresponding counts are three and nine.  The tensor rule itself is dimension-independent. \lecturetimestamp{00:42:55}}
\end{classroomqa}

\begin{classroomqa}
\audiencequestion{Can covariant and contravariant objects be described more geometrically?}
\lecturerresponse{Yes, they have a geometric interpretation, but pursuing it now would postpone electrodynamics even further.  For this course the transformation and index-manipulation rules are sufficient; the geometric distinction will become unavoidable in general relativity. \lecturetimestamp{00:43:51}}
\end{classroomqa}

\section{Symmetry, Antisymmetry, and the Electromagnetic Tensor}

For a generic rank-two tensor, exchanging the indices changes the component:
\begin{equation}
T^{\mu\nu}\neq T^{\nu\mu}.
\end{equation}
Indeed, \(A^\mu B^\nu\) and \(A^\nu B^\mu\) are different in general; the
specific components \(A^0B^1\) and \(A^1B^0\) need not agree.

There are two important special cases.  A symmetric tensor satisfies
\begin{equation}
T^{\mu\nu}=T^{\nu\mu},
\end{equation}
and a concrete symmetric combination is
\begin{equation}
A^\mu B^\nu+A^\nu B^\mu.
\end{equation}
An antisymmetric tensor satisfies
\begin{equation}
F^{\mu\nu}=-F^{\nu\mu},
\label{eq:sr6-antisymmetry}
\end{equation}
as does the combination
\begin{equation}
A^\mu B^\nu-A^\nu B^\mu.
\end{equation}
Setting the indices equal in \eqref{eq:sr6-antisymmetry} gives
\begin{equation}
F^{00}=-F^{00},
\qquad
F^{00}=0,
\end{equation}
and the same argument removes every diagonal component.  A \(4\times4\)
antisymmetric tensor therefore has six independent components.

\begin{figure}[htbp]
\centering
\includegraphics[width=0.90\textwidth]{lecture_06_figure_10.png}
\caption{Symmetric and antisymmetric rank-two tensors compared on the board.}
\lectureframe{00:47:20}
\end{figure}

Those six components are precisely enough for the three components of an
electric field and the three components of a magnetic field.  With the
sign convention written on the board,
\begin{equation}
F^{\mu\nu}
=
\begin{pmatrix}
0&-E_1&-E_2&-E_3\\
E_1&0&-B_3&B_2\\
E_2&B_3&0&-B_1\\
E_3&-B_2&B_1&0
\end{pmatrix}.
\label{eq:sr6-em-field-tensor}
\end{equation}
Reflection across the diagonal reverses every sign.  Electric and magnetic
fields are therefore not independent relativistic objects.  A boost can mix
them just as it mixes space and time: a field that is purely electric in one
frame may have a magnetic component in another.

\begin{figure}[htbp]
\centering
\includegraphics[width=0.90\textwidth]{lecture_06_figure_11.png}
\caption{The completed antisymmetric array whose six independent entries are \(\mathbf E\) and \(\mathbf B\).}
\lectureframe{00:50:20}
\end{figure}

The board carrying \eqref{eq:sr6-em-field-tensor} remains visible while the
discussion moves to particle dynamics.  We now reuse the symbol \(A_\mu\),
not for the generic vector of the preceding examples, but for the
electromagnetic four-potential.

\section{The Charged-Particle Action}

There are two possible starting points.  We could study the dynamics of the
field itself, which will lead to Maxwell's equations, or study a charged
particle moving in a prescribed field, which leads to the Lorentz force.
We begin with the particle.

At low velocity the target is
\begin{equation}
m\mathbf a
=
e\left(\mathbf E+\mathbf v\times\mathbf B\right),
\label{eq:sr6-nonrel-lorentz-force}
\end{equation}
with a factor \(1/c\) multiplying the magnetic term in the Gaussian-unit
version.  We set \(c=1\).  The two terms in
\eqref{eq:sr6-nonrel-lorentz-force} look different in three-vector notation,
but relativity should reveal them as parts of one object.

\begin{figure}[htbp]
\centering
\includegraphics[width=0.90\textwidth]{lecture_06_figure_12.png}
\caption{The nonrelativistic Lorentz force, the result the invariant action must reproduce.}
\lectureframe{00:54:00}
\end{figure}

\subsection{Free motion and the vector potential}

For a free relativistic particle,
\begin{equation}
S_{\mathrm{free}}
=
-m\int_1^2 d\tau
=
\int_{t_1}^{t_2}dt\,
\left[-m\sqrt{1-\dot{\mathbf x}^{\,2}}\right],
\label{eq:sr6-free-action}
\end{equation}
where
\begin{equation}
\dot{\mathbf x}^{\,2}
=
(\dot x^1)^2+(\dot x^2)^2+(\dot x^3)^2.
\end{equation}
The quantity in square brackets is the free-particle Lagrangian.  Its
momentum and energy were found earlier; we retain the action and ask what
field-dependent term may be added.

\begin{figure}[htbp]
\centering
\includegraphics[width=0.90\textwidth]{lecture_06_figure_13.png}
\caption{The free action as proper time accumulated along a worldline.}
\lectureframe{00:55:10}
\end{figure}

The basic electromagnetic variable is the four-vector potential
\(A_\mu(x)\).  The electric and magnetic fields will be derived from it.
Along a small worldline segment \(dx^\mu\), the contraction
\begin{equation}
A_\mu(x)\,dx^\mu
\end{equation}
is a scalar.  Adding those scalar increments from one endpoint to the other
gives a Lorentz-invariant interaction:
\begin{equation}
S_{\mathrm{int}}
=
-e\int_1^2 A_\mu(x)\,dx^\mu.
\label{eq:sr6-interaction-action}
\end{equation}
The coefficient \(e\) is the electric charge.  The overall sign belongs to
the historical charge convention; changing the convention changes the sign
assigned to \(e\), not the content of the theory.

\begin{figure}[htbp]
\centering
\includegraphics[width=0.90\textwidth]{lecture_06_figure_14.png}
\caption{A worldline segment contracted with the four-potential to form the scalar increment \(A_\mu dx^\mu\).}
\lectureframe{00:59:20}
\end{figure}

\subsection{Ordinary-time form}

To use the familiar Euler--Lagrange equations, divide and multiply
\eqref{eq:sr6-interaction-action} by \(dt\).  Since
\begin{equation}
\frac{dx^0}{dt}=1,
\qquad
\frac{dx^m}{dt}=\dot x^m,
\end{equation}
the complete action is
\begin{equation}
S
=
\int dt\,
\left[
-m\sqrt{1-\dot{\mathbf x}^{\,2}}
-eA_0(x,t)
-e\dot x^m A_m(x,t)
\right].
\label{eq:sr6-complete-particle-action}
\end{equation}
Thus
\begin{equation}
\boxed{
L
=
-m\sqrt{1-\dot{\mathbf x}^{\,2}}
-eA_0
-e\dot x^mA_m
}.
\label{eq:sr6-charged-lagrangian}
\end{equation}
For now \(A_\mu(x,t)\) is a known external field.  We vary the worldline,
not the field.

\begin{figure}[htbp]
\centering
\includegraphics[width=0.90\textwidth]{lecture_06_figure_15.png}
\caption{The invariant worldline coupling rewritten as an ordinary-time Lagrangian.}
\lectureframe{01:05:20}
\end{figure}

\section{Grinding Out the Lorentz Force}

For a spatial component \(x^m\), the Euler--Lagrange equation reads
\begin{equation}
\frac{d}{dt}\frac{\partial L}{\partial\dot x^m}
=
\frac{\partial L}{\partial x^m}.
\label{eq:sr6-EL}
\end{equation}
Differentiating with respect to velocity gives
\begin{equation}
\frac{\partial L}{\partial\dot x^m}
=
\frac{m\dot x_m}{\sqrt{1-\dot{\mathbf x}^{\,2}}}
-eA_m.
\label{eq:sr6-canonical-momentum}
\end{equation}
The letter \(m\) unfortunately serves both as the mass and as a spatial
index.  Context distinguishes them: the first \(m\) in
\eqref{eq:sr6-canonical-momentum} is mass; the subscript \(m\) labels one of
the three spatial components.

\begin{figure}[htbp]
\centering
\includegraphics[width=0.90\textwidth]{lecture_06_figure_16.png}
\caption{The momentum derivative contains both the mechanical term and the vector-potential term.}
\lectureframe{01:09:20}
\end{figure}

\begin{classroomqa}
\audiencequestion{Is the first attempt at the left-hand side missing a term?}
\lecturerresponse{Yes.  Differentiating only the free-particle term omits the velocity dependence of \(-e\dot x^mA_m\).  The full derivative is \(m\dot x_m/\sqrt{1-\dot{\mathbf x}^{\,2}}-eA_m\).  The correction is essential because its time derivative supplies part of both the electric and magnetic force. \lecturetimestamp{01:08:18}}
\end{classroomqa}

The right-hand side of \eqref{eq:sr6-EL} is
\begin{equation}
\frac{\partial L}{\partial x^m}
=
-e\,\partial_mA_0
-e\,\dot x^n\partial_mA_n.
\label{eq:sr6-position-derivative}
\end{equation}
Therefore
\begin{equation}
\frac{d}{dt}
\left(
\frac{m\dot x_m}{\sqrt{1-\dot{\mathbf x}^{\,2}}}
-eA_m
\right)
=
-e\,\partial_mA_0
-e\,\dot x^n\partial_mA_n.
\label{eq:sr6-EL-expanded}
\end{equation}

\begin{classroomqa}
\audiencequestion{Have covariant and contravariant spatial components been interchanged in this calculation?}
\lecturerresponse{They have been moved up and down, but for Latin spatial indices the Euclidean spatial block of the metric is the identity.  Thus \(x_m=x^m\) and \(A_m=A^m\).  Only a time index produces the sign change. \lecturetimestamp{01:09:29}}
\end{classroomqa}

\begin{classroomqa}
\audiencequestion{Are repeated Latin indices summed too?}
\lecturerresponse{Yes.  A repeated Latin index is summed from \(1\) to \(3\); a repeated Greek index is summed from \(0\) to \(3\).  A lone index such as the \(m\) labeling one spatial equation is free and is not summed.  The mass symbol is not an index. \lecturetimestamp{01:10:00}}
\end{classroomqa}

\begin{figure}[htbp]
\centering
\includegraphics[width=0.90\textwidth]{lecture_06_figure_17.png}
\caption{The Euler--Lagrange equation before the vector potential's total time derivative is expanded.}
\lectureframe{01:14:20}
\end{figure}

\subsection{The chain rule along the moving worldline}

Even a time-independent field changes along the particle's path because the
particle samples different positions.  The total derivative is
\begin{equation}
\frac{dA_m(x(t),t)}{dt}
=
\partial_0A_m+\dot x^n\partial_nA_m.
\label{eq:sr6-convective-derivative}
\end{equation}
The first term is explicit time dependence; the second is implicit time
dependence through \(x^n(t)\).

\begin{figure}[htbp]
\centering
\includegraphics[width=0.90\textwidth]{lecture_06_figure_18.png}
\caption{The vector potential changes explicitly with time and implicitly as the particle moves through space.}
\lectureframe{01:17:20}
\end{figure}

Substituting \eqref{eq:sr6-convective-derivative} into
\eqref{eq:sr6-EL-expanded} and moving the field terms to the right gives
\begin{equation}
\begin{aligned}
m\frac{d}{dt}
\left(
\frac{\dot x_m}{\sqrt{1-\dot{\mathbf x}^{\,2}}}
\right)
&=
e\left(\partial_0A_m-\partial_mA_0\right)
\\
&\quad +e\dot x^n
\left(\partial_nA_m-\partial_mA_n\right).
\end{aligned}
\label{eq:sr6-grouped-force}
\end{equation}
The calculation has finally found its organizing principle.  The first
bracket contains no velocity; the second is linear in velocity.  This is
exactly the split between electric and magnetic forces.

\begin{figure}[htbp]
\centering
\includegraphics[width=0.90\textwidth]{lecture_06_figure_19.png}
\caption{The force terms grouped into a velocity-independent electric part and a velocity-dependent antisymmetric part.}
\lectureframe{01:23:40}
\end{figure}

Define the electric components in the convention of the derivation by
\begin{equation}
E_m=\partial_0A_m-\partial_mA_0.
\label{eq:sr6-electric-from-potential}
\end{equation}
The spatial antisymmetric combination
\begin{equation}
C_{mn}
\equiv
\partial_nA_m-\partial_mA_n
\label{eq:sr6-spatial-curl-tensor}
\end{equation}
contains the three independent components of the curl of the vector
potential.  Its orientation is chosen so that
\begin{equation}
\dot x^nC_{mn}
=
\left(\mathbf v\times\mathbf B\right)_m.
\end{equation}
Equation \eqref{eq:sr6-grouped-force} is therefore the relativistic momentum
form of \eqref{eq:sr6-nonrel-lorentz-force}.  The factor under the time
derivative is the relativistic momentum per unit mass, not the Newtonian
velocity alone.

The derivation is worth repeating by hand: begin with
\eqref{eq:sr6-charged-lagrangian}, compute both sides of
\eqref{eq:sr6-EL}, apply the chain rule, and group terms according to their
velocity dependence.  The exercise is simultaneously practice in
variational mechanics, multivariable calculus, and electrodynamics.

\begin{classroomqa}
\audiencequestion{Why choose \(-e\int A_\mu dx^\mu\) as the electromagnetic interaction?}
\lecturerresponse{Experiment ultimately decides the coupling, but two structural requirements make this the natural first candidate.  The action is assembled from local increments along the trajectory, and each increment must be a scalar.  The scalar \(dx^\mu dx_\mu=d\tau^2\) is already represented by the free action.  Contracting \(dx^\mu\) with the field \(A_\mu\) is the simplest new field-dependent scalar.  More elaborate scalars can be written, and gauge invariance will later sharpen the choice. \lecturetimestamp{01:25:34}}
\end{classroomqa}

\begin{classroomqa}
\audiencequestion{Can Maxwell's equations also be derived from least action?}
\lecturerresponse{Yes.  The present action governs a particle in a given electromagnetic field.  To obtain Maxwell's equations we must add the field's own action and vary the field itself, just as we did for the scalar field. \lecturetimestamp{01:28:58}}
\end{classroomqa}

\begin{classroomqa}
\audiencequestion{Where is the summation in the grouped force equation?}
\lecturerresponse{The free \(m\) labels one of the three component equations and is not summed.  The repeated \(n\) in \(\dot x^n(\partial_nA_m-\partial_mA_n)\) is summed over \(1,2,3\).  Written without Einstein notation, that term would carry an explicit \(\sum_{n=1}^{3}\). \lecturetimestamp{01:29:26}}
\end{classroomqa}

\begin{classroomqa}
\audiencequestion{Why identify \(\partial_0A_m-\partial_mA_0\) as the electric field rather than assume the usual formula in advance?}
\lecturerresponse{The action has produced a force term that depends on position and time but not on the particle velocity.  That is precisely the operational role of the electric field.  The remaining antisymmetric terms multiply the velocity and therefore have the magnetic role.  One may regard this calculation as discovering the fields from the force law rather than inserting separate definitions beforehand. \lecturetimestamp{01:31:05}}
\end{classroomqa}

\section{The Manifestly Covariant Force Law}

The left-hand side of \eqref{eq:sr6-grouped-force} is not yet written in an
obviously four-dimensional form.  Introduce the four-velocity
\begin{equation}
u^\mu=\frac{dx^\mu}{d\tau}.
\label{eq:sr6-four-velocity}
\end{equation}
Its spatial components obey
\begin{equation}
u^m
=
\frac{dx^m}{d\tau}
=
\frac{\dot x^m}{\sqrt{1-\dot{\mathbf x}^{\,2}}}.
\end{equation}
Multiplying the spatial equation by \(dt/d\tau\) converts
\(d/dt\) into \(d/d\tau\):
\begin{equation}
\frac{dt}{d\tau}\frac{d}{dt}
=
\frac{d}{d\tau}.
\end{equation}
Hence the left-hand side becomes the spatial part of the proper
acceleration,
\begin{equation}
m\frac{d^2x^m}{d\tau^2}.
\end{equation}

The right-hand side can be written without separating time and space:
\begin{equation}
m\frac{d^2x^m}{d\tau^2}
=
e\left(\partial_\nu A_m-\partial_mA_\nu\right)
\frac{dx^\nu}{d\tau}.
\label{eq:sr6-spatial-covariant-force}
\end{equation}
For \(\nu=0\) this reproduces the electric bracket; for
\(\nu=1,2,3\) it reproduces the magnetic bracket.

Both sides now have the transformation character of a four-vector.  The
three spatial equations are therefore components of
\begin{equation}
\boxed{
m\frac{d^2x^\mu}{d\tau^2}
=
e\left(\partial_\nu A^\mu-\partial^\mu A_\nu\right)u^\nu
}.
\label{eq:sr6-covariant-force}
\end{equation}
If one defines
\begin{equation}
F_{\mu\nu}
\equiv
\partial_\mu A_\nu-\partial_\nu A_\mu,
\label{eq:sr6-F-from-A}
\end{equation}
then the index order and interaction sign used in this lecture give the
equivalent compact form
\begin{equation}
m\frac{du^\mu}{d\tau}
=
-eF^\mu{}_\nu u^\nu.
\label{eq:sr6-covariant-force-F}
\end{equation}
\editorialnote{Electrodynamics texts distribute signs differently among the metric signature, the definition of \(A_\mu\), the charge, and \(F_{\mu\nu}\).  Equations \eqref{eq:sr6-covariant-force} and \eqref{eq:sr6-covariant-force-F} retain the sign and index order obtained from the action \(-e\int A_\mu dx^\mu\); the invariant content is the antisymmetrized derivative of the potential.}

\begin{figure}[htbp]
\centering
\includegraphics[width=0.90\textwidth]{lecture_06_figure_20.png}
\caption{The three spatial equations assembled into a four-vector equation using proper time.}
\lectureframe{01:42:40}
\end{figure}

\begin{classroomqa}
\audiencequestion{Why is it legitimate to add a fourth, time-component equation?  Is \(d^2x^0/d\tau^2\) simply zero?}
\lecturerresponse{It is not generally zero: \(x^0=t\), but \(dt/d\tau=\gamma\) changes when the particle's speed changes.  The decisive argument is covariance.  The scalar action guarantees Lorentz-invariant equations; once the spatial components are recognized as components of one four-vector equation, Lorentz transformations mix them with the time component.  The fourth equation must accompany the other three. \lecturetimestamp{01:42:40}}
\end{classroomqa}

\subsection{What the fourth component means}

The three spatial components are the relativistic statement
\begin{equation}
\frac{d\mathbf p}{dt}=\mathbf F.
\end{equation}
Energy and momentum are themselves the components of a four-vector, so the
time component must describe the corresponding energy balance.  In
three-vector language,
\begin{equation}
\frac{dK}{dt}=\mathbf F\cdot\mathbf v.
\label{eq:sr6-work-energy}
\end{equation}
It records the rate at which the field does work on the particle.  The
magnetic contribution is perpendicular to \(\mathbf v\) and does no work;
the electric contribution changes the particle energy.

\begin{classroomqa}
\audiencequestion{If covariance supplies the fourth equation formally, how do we know it has physical meaning?}
\lecturerresponse{The first three components say that force changes momentum.  The fourth says that work changes energy: \(dK/dt=\mathbf F\cdot\mathbf v\).  Since energy and momentum form one four-vector, their response to the field must also be one four-component statement. \lecturetimestamp{01:45:20}}
\end{classroomqa}

\begin{figure}[htbp]
\centering
\includegraphics[width=0.90\textwidth]{lecture_06_figure_21.png}
\caption{The spatial law \(d\mathbf p/dt=\mathbf F\) and its time-component partner, the work--energy relation.}
\lectureframe{01:47:10}
\end{figure}

The method is more important than the final compression.  We required the
action to be a scalar before varying it.  The resulting equations could
then be reorganized into tensors without repairing their symmetry after the
fact.  In modern field theory this becomes a general design rule: construct
the Lagrangian to respect the symmetries of the problem.

\section{Locality, Lorentz Invariance, and the Next Symmetry}

Three principles recur throughout relativistic field theory.  Two have
already been used; the third will be developed next.

\subsection{Locality}

Locality says that the change of a particle or field is controlled by what
happens infinitesimally nearby in spacetime.  A particle responds to the
field evaluated on its worldline, not to an unrelated value far away.  A
field at one point responds to the field and its derivatives in a
neighborhood of that point, not directly to arbitrary remote values or to
future data.

For a field, the action therefore takes the local form
\begin{equation}
S
=
\int d^3x\,dt\,
\mathcal L\!\left(\phi,\partial_\mu\phi\right).
\label{eq:sr6-local-action}
\end{equation}
The integral adds one contribution from each spacetime point.  First
derivatives encode comparison with neighboring points, and the variational
principle produces differential equations.  One could imagine nonlocal
products coupling fields at widely separated points, but their equations
would make the response here depend directly on conditions there.  That is
not the local structure being imposed.  Actions with higher derivatives are
possible and form an interesting subject, but they are not needed for the
present construction.

\begin{figure}[htbp]
\centering
\includegraphics[width=0.90\textwidth]{lecture_06_figure_22.png}
\caption{Locality expressed as a spacetime integral of a density built from fields and their nearby derivatives.}
\lectureframe{01:51:00}
\end{figure}

\begin{classroomqa}
\audiencequestion{Does \(\partial_\mu\phi\) represent four derivative components?}
\lecturerresponse{Yes.  It denotes the time derivative and the three spatial derivatives of \(\phi\), collected as one covariant four-vector. \lecturetimestamp{01:52:25}}
\end{classroomqa}

\begin{classroomqa}
\audiencequestion{How is \emph{nearby} quantified in the locality principle?}
\lecturerresponse{By derivatives.  A first derivative compares the field at a point with its value at infinitesimally neighboring points.  A local Lagrangian depends on the fields and a finite set of such derivatives at the same spacetime point. \lecturetimestamp{01:52:55}}
\end{classroomqa}

\begin{classroomqa}
\audiencequestion{How does the field version of locality translate to the particle calculation just completed?}
\lecturerresponse{The particle Lagrangian is local along the worldline: it depends on neighboring worldline positions through the velocity and on the electromagnetic potential evaluated at the particle's current position.  The particle's acceleration here does not depend directly on the field at a distant point. \lecturetimestamp{01:53:31}}
\end{classroomqa}

\subsection{Lorentz invariance}

Lorentz invariance is implemented by building the Lagrangian from scalars.
Observers may decompose those scalars into different time and space
components, but they agree on the action.  The tensor equations obtained by
variation consequently have the same form in every inertial frame.

\subsection{Gauge invariance}

Gauge invariance is the remaining principle.  It is announced but not yet
defined here.  In the next step it will constrain which features of the
four-potential are physically meaningful and will guide the field action
that produces Maxwell's equations.

\begin{figure}[htbp]
\centering
\includegraphics[width=0.90\textwidth]{lecture_06_figure_23.png}
\caption{The three organizing principles: locality, Lorentz invariance, and gauge invariance.}
\lectureframe{01:55:40}
\end{figure}

These principles recur in general relativity, quantum electrodynamics,
Yang--Mills theory, and the Standard Model.  Details such as the number of
particle species vary from theory to theory; locality and symmetry provide
the deeper framework.  The practical assignment is therefore to retrace
the Euler--Lagrange calculation until every index and every grouped term can
be accounted for.  Gauge invariance will then enter as the next constraint
on the same electromagnetic action.
